\section{Analytic link smearing}
\label{sec:analytic}

A method of smearing link variables which is analytic, and hence
differentiable, everywhere in the finite complex plane
can be defined as follows.
Let $C_\mu(x)$ denote the weighted sum of the perpendicular staples which
begin at lattice site $x$ and terminate at neighboring site 
$x\!+\!\hat{\mu}$:
\begin{eqnarray}
 C_\mu(x)&=&\sum_{\nu\neq \mu}\rho_{\mu\nu}\biggl(
 U_\nu(x) U_\mu(x\!+\!\hat{\nu}) U_\nu^\dagger(x\!+\!\hat{\mu})\nonumber\\
&&+ U^\dagger_\nu(x\!-\!\hat{\nu}) U_\mu(x\!-\!\hat{\nu})
  U_\nu(x\!-\!\hat{\nu}\!+\!\hat{\mu})
\biggr), \label{eq:Cdef}
\end{eqnarray}
where $\hat{\mu},\hat{\nu}$ are vectors in directions $\mu,\nu$,
respectively, having the length
of one lattice spacing in that direction (fundamental lattice
translation vectors).  The weights $\rho_{\mu\nu}$ are tunable real
parameters. Then the matrix $Q_\mu(x)$, defined in $SU(N)$ by
\begin{eqnarray}
Q_\mu(x) &=& 
  \frac{i}{2}\Bigl(\!\Omega^\dagger_\mu(x)\!-\!\Omega_\mu(x)\!\Bigr)
   \!-\!\frac{i}{2N}\Tr\Bigl(\!\Omega^\dagger_\mu(x)
  \!-\!\Omega_\mu(x)\!\Bigr),\nonumber\\
\Omega_\mu(x) &=& C_\mu(x)\ U_\mu^\dagger(x),
   \quad\mbox{(no summation over $\mu$)} \label{eq:Qdef}
\end{eqnarray}
is Hermitian and traceless, and hence, $e^{i Q_\mu(x)}$ is an
element of $SU(N)$.  
We use this fact to define an iterative, analytic link
smearing algorithm in which the links $U_\mu^{(n)}(x)$ at step $n$
are mapped into links $U_\mu^{(n+1)}(x)$ using
\begin{equation}
U^{(n+1)}_\mu(x)=\exp\Bigl(i Q_\mu^{(n)}(x)\Bigr)\ U_\mu^{(n)}(x).
\label{eq:Ustout}
\end{equation}
The fact that $e^{iQ_\mu(x)}$
is an element of $SU(N)$ guarantees that $U^{(n+1)}_\mu(x)$ is also an element
of $SU(N)$, eliminating the need for a projection back onto the gauge group.
This fuzzing step can be iterated $n_\rho$ times to finally
produce link variables which we call {\em stout links}\cite{pub},
denoted by $\tilde{U}_\mu(x)$:
\begin{equation}
 U\rightarrow U^{(1)}\rightarrow U^{(2)}\rightarrow\cdots\rightarrow
  U^{(n_\rho)}\equiv \tilde{U}.
\end{equation}

One common choice of the staple weights is
\begin{equation}
   \rho_{jk}=\rho,\qquad \rho_{4\mu}=\rho_{\mu 4}=0,
\end{equation}
which yields a three-dimensional scheme in which only the spatial
links are smeared.  Another common choice is an isotropic four-dimensional
scheme in which all weights are chosen to be the same
$\rho_{\mu\nu}=\rho$.  Note that such a scheme 
was used in Ref.~\cite{luscher} as a field transformation
to eliminate a certain interaction term from the most general $O(a^2)$
on-shell improved gauge action.

It is not difficult to show that, given an appropriate choice of weights
$\rho_{\mu\nu}$, the stout links have symmetry transformation
properties identical to those of the original link variables.
Under any local gauge transformation $G(x)$,
the link variables transform according to 
$U_\mu(x)\rightarrow G(x)\ U_\mu(x)\ G^\dagger(x\!+\!\hat{\mu})$,
where the $G(x)$ are $SU(N)$ matrices.  Thus,
under such a gauge transformation,
\begin{eqnarray}
Q_\mu(x)&\rightarrow& G(x)\ Q_\mu(x)\ G^\dagger(x),\\
e^{iQ_\mu(x)}&\rightarrow& G(x)\ e^{iQ_\mu(x)}\ G^\dagger(x),\\
U^{(1)}_\mu(x)&\rightarrow& G(x)\ e^{iQ_\mu(x)} G^\dagger(x)
 G(x)U_\mu(x)\ G^\dagger(x\!+\!\hat{\mu})\nonumber\\
 &=& G(x)\ U^{(1)}_\mu(x)\ G^\dagger(x\!+\!\hat{\mu}),
\end{eqnarray}
as required.  One can also easily show that if the weights $\rho_{\mu\nu}$
respect the rotation and reflection symmetries of the lattice,
the stout links obey the
required transformation properties under all rotations and all
reflections in a plane containing the link.  Lastly, consider a
reflection in a plane normal to a link $U_\mu(x)$ and passing through its
midpoint $x\!+\!\frac{1}{2}\hat{\mu}$. Under such an operation, the
link $U_\mu(x)$ transforms according to 
$U_\mu(x)\rightarrow U_{-\mu}(x\!+\!\hat{\mu})=U_\mu^\dagger(x)$ and
for $\nu\neq \mu$, $U_\nu(x)\leftrightarrow U_\nu(x\!+\!\hat{\mu})$ and
$U_\nu^\dagger(x\!+\!\hat{\mu})\leftrightarrow U_\nu^\dagger(x)$, yielding
$C_\mu(x)\rightarrow C_\mu^\dagger(x)$ assuming the $\rho_{\mu\nu}$
are real, so that
\begin{eqnarray}
 \Omega_\mu(x) &\rightarrow & C_\mu^\dagger(x)U_\mu(x)=U_\mu^\dagger(x)
 \Omega_\mu^\dagger(x) U_\mu(x),\\
 Q_\mu(x)&\rightarrow & -U_\mu^\dagger(x)Q_\mu(x)U_\mu(x),\\
 U^{(1)}_\mu(x)&\rightarrow & U_\mu^\dagger(x) e^{-iQ_\mu(x)}
   = U_\mu^{(1)\dagger}(x),
\end{eqnarray}
as required.  Given that $U^{(1)}_\mu(x)$ transforms under all symmetry
operations in the same manner as $U_\mu(x)$, it follows that the stout
links $\tilde{U}_\mu(x)$ have symmetry transformation properties identical
to those of the original link variables.

\begin{figure}
\includegraphics[width=0.98\textwidth]{images/paths}
\caption[figF]{The expansion up to first order in the $\rho_{\mu\nu}$
 of the new link variable
 $U^{(1)}$ in terms of paths of the original links.  Each closed loop
 includes a trace with a factor $1/N$ in $SU(N)$.
\label{fig:paths}}
\end{figure}

Since the exponential function has a power series expansion with
an infinite radius of convergence, each stout link may be viewed as
an incredibly large and complicated sum of paths.  For small $\rho_{\mu\nu}$,
the paths which make up the link variable $U^{(1)}$ to first order
in the $\rho_{\mu\nu}$ are shown in Fig.~\ref{fig:paths}.  Note that the 
standard smearing method, defined by
\begin{equation} 
U^{(1)}_\mu(x) = {\cal P}_{SU(3)}\ \biggl\{ 
 U_\mu(x) + C_\mu(x) \biggr\},
\label{eq:Ufuzz}
\end{equation}
where ${\cal P}_{SU(3)}$ denotes the projection into $SU(3)$,
yields the same sum of paths at first order in $\rho_{\mu\nu}$.

A few further remarks are worthy of note.  First, an alternative smearing
scheme in which the weights $\rho_{\mu\nu}$ are chosen to be imaginary
does not reduce the couplings to the high-frequency modes of the theory.
Secondly, it is not possible to remove the leading $O(a^2)$ discretization
errors in the Wilson gauge action by expressing the action in terms
of stout links; diminishing the $O(\alpha_s a^2)$ errors by reducing their
tadpole contributions is the best which can be achieved.  Thirdly, the
use of stout links in Symanzik-improved actions may eliminate the need
to use tadpole-improvement in determining the couplings in the action.
Lastly, the definition of $C_\mu(x)$ given in Eq.~(\ref{eq:Cdef}) is not
unique. Analytic smearing schemes can be constructed as outlined above
using sums of other paths from site $x$ to $x\!+\!\hat{\mu}$, such as
those used in fat link\cite{fat1} and HYP\cite{hyp} smearings.

